\documentclass[11pt,a4paper]{article}
\usepackage[T1]{fontenc}
\usepackage{lmodern,amsmath,amssymb,amsthm,mathtools,booktabs,microtype}
\usepackage[margin=26mm]{geometry}
\usepackage[colorlinks=true,linkcolor=blue,citecolor=blue,urlcolor=blue]{hyperref}
\setlength{\emergencystretch}{2em}
\newtheorem{theorem}{Theorem}[section]
\newtheorem{proposition}[theorem]{Proposition}
\newtheorem{lemma}[theorem]{Lemma}
\newtheorem{corollary}[theorem]{Corollary}
\theoremstyle{definition}\newtheorem{definition}[theorem]{Definition}
\theoremstyle{remark}\newtheorem{remark}[theorem]{Remark}
\newcommand{\Q}{\mathbb Q}\newcommand{\Z}{\mathbb Z}
\newcommand{\R}{\mathbb R}\newcommand{\A}{\mathbb A}
\newcommand{\PP}{\mathbb P}\newcommand{\sqcl}[1]{[#1]}
\DeclareMathOperator{\Tr}{Tr}\DeclareMathOperator{\Nm}{N}
\DeclareMathOperator{\N}{N}\DeclareMathOperator{\rad}{rad}
\DeclareMathOperator{\sfpart}{sf}
\title{Fresh-Prime and Fixed-Scale Obstructions to Joint Pole Realization}
\author{}\date{16 November 2026}
\begin{document}\maketitle
\begin{abstract}
The selected-prime completeness family for two six-contact detectors
forces their nonzero pole inputs to have different rational square
classes. Requiring equal first classes instead gives a genus-four
curve at each rational scale greater than four. For disjoint
squarefree contact polynomials of a common positive even degree $m$,
we compute the genus of the simultaneous square-class correspondence
as $1+2^{r-3}(rm-4)$ for $r$ poles. The ten-contact cases have genera
nine and twenty-seven for two and three poles. These obstructions
apply to the specified witness family and fixed-scale slope
arguments, not to all possible Diophantine constructions.
\end{abstract}

\section{Scope and square classes}
For $b\in\Q^\times$, let $\sqcl b$ denote its class in
$\Q^\times/\Q^{\times2}$. Exact transfer of pole witnesses can
reduce the retained parameters when both the inputs and their
specified complements have matching square classes~\cite{JointRealization}.
The results here concern two constraints on that compatibility:
the selected-prime construction forces different first classes,
and at fixed scales equality of those classes defines curves of
positive genus. Finiteness at fixed scales restricts an argument
requiring arbitrarily many slopes; it does not exclude a suitable
finite choice for each arithmetic input.

\section{Selected prime classes in the six-contact family}
The six-contact conditional construction uses
\begin{equation}\label{eq:F}
 \Theta=\left\{0,-\tfrac12,-\tfrac14,-\tfrac25,-1,-\tfrac13\right\},
 \quad F(t)=\prod_{\theta\in\Theta}(t-\theta),\quad
 t_0=\lambda,\quad t_1=a\lambda.
\end{equation}
Here $a=\beta^d+T^d>4$, with the notation and hypotheses of
\cite{Input13,UnitNorm}. Write $W_j=(W_j)_0+(W_j)_1\sqrt2$
and $\N(W_j)=(W_j)_0^2-2(W_j)_1^2$. On the source locus,
$t_j\notin\Theta\cup\{1\}$, $W_j\ne0$, and $(W_j)_0\ne0$.
The nonzero pole inputs have the form
\begin{equation}\label{eq:actualclass}
 b_j=F(t_j)\left(\frac{\N(W_j)}{2(t_j-1)^4(W_j)_0}\right)^2.
\end{equation}
Thus $\sqcl {b_j}=\sqcl {F(t_j)}$. This statement is about the literal
source equations, not a numerical approximation to them.

\begin{proposition}[Incompatibility of the selected prime classes]
\label{prop:fresh}
In the simultaneous-prime completeness construction for
\eqref{eq:F}, suppose each distinct contact coefficient has its own
selected prime, with valuation one in its incident contact numerator,
and with no incidence in any other distinct coefficient. These are the
retained fresh-prime properties of that construction. Then
$b_1/b_0$ in \eqref{eq:actualclass} is not a rational square.
\end{proposition}
\begin{proof}
The nonzero numbers in $\Theta$ have absolute values between $1/4$ and~1.
If $a>4$, the nonzero elements of $\Theta/a$ have absolute value less
than $1/4$. Hence the contact coefficient sets $\Theta$ and $\Theta/a$
intersect only at~0. Choose a selected prime belonging to a nonzero
contact of the first detector. It has valuation one in $F(t_0)$ and
valuation zero in $F(t_1)$; the fixed scale and denominator primes are
excluded in the retained prime construction. The ratio in
\eqref{eq:actualclass} therefore has odd valuation at this prime.
\end{proof}

The selected-prime hypotheses therefore exclude equality of the
two input square classes. This concerns the stated completeness
family; it does not prove that every possible witness has different
classes or give a lower bound on the number of variables.

\begin{proposition}[The fixed-scale curve]\label{prop:genus}
Fix a rational $a>4$. Write $F(t)=tG(t)$, with $G$ the monic degree-five
polynomial defined by \eqref{eq:F}. For nonzero $\lambda$ away from
all contacts,
\[
 \sqcl {F(\lambda)}=\sqcl {F(a\lambda)}
 \quad\Longleftrightarrow\quad
 \exists y\in\Q:\ y^2=aG(\lambda)G(a\lambda).
\]
The smooth projective model of this curve has genus four. In particular,
for each fixed rational $a>4$ it has only finitely many rational
$\lambda$, by Faltings' theorem.
\end{proposition}
\begin{proof}
Two nonzero rationals have the same square class exactly when their
product is a square, and
$F(\lambda)F(a\lambda)=a\lambda^2G(\lambda)G(a\lambda)$.
The ten roots are the nonzero contacts and their quotients by~$a$.
They are distinct by the argument in Proposition~\ref{prop:fresh}.
A squarefree polynomial of degree ten defines a double cover of
$\mathbf P^1$ branched at ten points and unramified at infinity.
Riemann--Hurwitz gives $2g-2=-4+10$, so $g=4$.
Faltings' theorem~\cite{Faltings} applies to this smooth projective
curve over $\Q$. The more general contact-parity covers, including
their fixed-scale finiteness statement, are treated in~\cite{ContactCovers}.
\end{proof}

\begin{remark}
The assertion is pointwise in $a$, not a uniform bound on rational
points as $a$ varies. It does not rule out one suitable finite witness
for every required input. It does rule out proving the desired match by
an arbitrarily-large-slope argument for a fixed scale. The points at
infinity do not give finite witnesses to the source equations.
\end{remark}

\subsection{The ten-contact and three-pole variants}
A similar calculation applies to the twenty-five-variable family.
At fixed values of its other retained coordinates, and with
nonzero rational scales $a_j$, write
\[
 \Xi=\{0,-1/2,-1,1/2399,-32/407\},\qquad
 \Xi_j=(4j+2+\Xi)\cup(4j+4+\Xi),
\]
\[
 H_j(\lambda)=\prod_{\xi\in\Xi_j}(a_j\lambda-\xi),\qquad a_0=1.
\]
Each $\Xi_j$ has ten distinct nonzero elements. The literal pole
input is $\rho H_j$, with $\rho\in\Q^\times$, times a nonzero rational square, so equality of
first square classes is equivalent to equality of the classes of
$H_j$. The common factor $\rho$ cancels.

\begin{proposition}\label{prop:multigenus}
Let $r\ge2$, and let $H_0,\ldots,H_{r-1}$ be squarefree polynomials
of the same positive even degree $m$, with disjoint root sets over
$\overline\Q$. The smooth projective model of the simultaneous
square-class correspondence
\[
 y_i^2=H_0(\lambda)H_i(\lambda)\quad(1\le i<r)
\]
is geometrically connected and has genus
\[
                         g=1+2^{r-3}(rm-4).
\]
For the ten-contact family, outside the finite set of contact-ratio
collisions this gives genus $9$ for two poles and genus $27$ for three.
\end{proposition}
\begin{proof}
The $r-1$ functions $H_0H_i$ are independent modulo squares in
$\overline\Q(\lambda)^\times$: a root unique to $H_i$ detects its
exponent in any putative square product. The function-field extension
therefore has degree $d=2^{r-1}$ and is geometrically connected.
At each of the $rm$ distinct roots, the nonzero parity vector gives
inertia of order two. Thus each branch point contributes $d/2$ to
Riemann--Hurwitz. There is no ramification at infinity because each
$H_0H_i$ has even degree. Hence
\[
             2g-2=-2d+rm\,d/2,
\]
which is the displayed formula. In the ten-contact application,
root sets meet exactly when $a_j/a_i=\xi_j/\xi_i$ for some contacts;
this is a finite exceptional collection of ratios. Taking $m=10$
and $r=2,3$ gives the asserted genera.
\end{proof}

The statement freezes the scales and all other retained variables
before varying the slope. It does not assert finiteness in the whole
parameter space. The additional complement-class requirements for joint pole
realization~\cite{JointRealization} only restrict this correspondence further.
As in Proposition~\ref{prop:genus}, the finiteness conclusion is an
obstruction to retaining the old unlimited-slope proof, not a proof
that no finite choice can work.


\section{Exact verification}
The program
\nolinkurl{computation/pole_parameters/joint_realization_v29/check.py}
verifies the fixed-scale contact-root computation and authenticates
the source files used for the pole-class identity. Its other checks
concern the transfer identities and rational examples
in~\cite{JointRealization}. The branch calculation above proves the
genus formula for all the stated inputs; rational-point finiteness
uses Faltings' theorem, rather than a finite search.

\begin{thebibliography}{99}
\bibitem{Input13} \emph{Thirteen Variables from Direct Elliptic Index Comparison},
project source manuscript,
\nolinkurl{papers/conditional_few_variables/direct_input13_20261110_v22.tex}.

\bibitem{UnitNorm} \emph{A single norm determinant for the 25- and 13-variable constructions},
project manuscript, 14 November 2026,
\nolinkurl{papers/finite_algebra/unit_norm_20261114_v28.tex}.

\bibitem{Faltings} G.~Faltings, \emph{Endlichkeitss\"atze f\"ur abelsche
Variet\"aten \"uber Zahlk\"orpern}, Invent. Math. \textbf{73} (1983),
349--366, Satz~7, p.~365; erratum, \textbf{75} (1984), 381.

\bibitem{JointRealization}
\emph{Joint realization of compatible pole conditions},
project manuscript, 14 November 2026,
\nolinkurl{papers/pole_parameters/joint_pole_realization_20261114_v29.tex}.

\bibitem{ContactCovers}
\emph{Quadratic Covers of Two Contact Functions},
project manuscript, 10 November 2026,
\nolinkurl{negative_results/papers/pole_parameters/quadratic_contact_covers_20261110_v21.tex}.
\end{thebibliography}
\end{document}
